\chapter{KLRW Category}

\section{Overview}
\section{Computable Add}


\begin{definition}
 Let $\C$ be a preadditive category. Then the \emph{computable additive completion}
 of $\C$, denoted $\operatorname{CMat}(\C)$, has
 \begin{itemize}
  \item as objects, lists of elements in $\C$.
  \item as morphisms, dependent matrices of morphisms in $\C$.
  Specifically, if $(A_0, A_0, \ldots, A_n), (B_0, B_1, \ldots, B_m) \in \Ob\big((\Mat (\C))\big)$,
  then $\Hom\big((A_0, A_1, \ldots, A_n), (B_0, B_1, \ldots, B_n)\big) =
  \{\begin{bmatrix}
    a_{00} & \cdots & a_{0a} \\ \cdots & \cdots & \cdots \\ a_{m0} & \cdots & a_{ma}
  \end{bmatrix}: a_{ij}\in \Hom(A_j, B_i), i\in \{0, \ldots, m-1\}, j\in \{0, \ldots, n-1\}\}
  $
  \item composition is defined by matrix multiplication
  \item identities are given by the identity matrix
  $$ (\idf_{(A_0, \ldots, A_m)})_{i, j} = \begin{cases}
  \idf_{A_i} & \text{reinterpreted as an element of }\Hom(A_i, A_j)
  \text{ by casting along the equality if }i = j
  \\ 0 & \text{ if }i\ne j
  \end{cases}$$.
 \end{itemize}
\end{definition}
  There is a fully faithful functor $G: \C \Rightarrow \CMat(\C)$.

  There is a fully faithful embedding functor $F: \CMat(\C) \Rightarrow \Mat (\C)$.
  It is (nonconstructively) essentially surjective, so $F$ is an equivalence
  of categories.
\begin{definition}
  Given $A = (A_0, \ldots, A_n), B=(B_0, \ldots, B_n)\in \CMat(\C)$, the
  \emph{computable biproduct} of $A$ and $B$ is defined by $A\boxplus_k B = (A_0, \ldots, A_n, B_0, \ldots, B_m)$
\end{definition}
\begin{definition}
  For any $Z\in \CMat (\C)$, the \emph{index type}, denoted $\iota_Z$, is the type
  $\{ 0, \ldots, (\len(Z)-1)\}$, and the indexing function, $X_2: \iota_Z\to \C$,
  is defined by $X_Z(i) = Z_i$.
\end{definition}

\begin{definition}
  Let $n\in \N$. A \emph{positioning} of $1$ black strand with $n$ red strands is
  an element of $\{0, \ldots, n+1\}$. Given positionings of $X$
  and $Y$, the strand set $StrandSpace_{X, Y} = \N$, representing
  the number of dots. (Note here that $\StrandSpace_{X, Y}$ does not depend on
  $X$ or $Y$, but with more black strands it might).
\end{definition}

\begin{definition}
  If $X, Y, Z$ are positionings of $1$ black strand and $n$ red strands
  and $a\in \StrandSpace_{X, Y}$, $b\in \StrandSpace_{Y, Z}$, then their
  composition is given by
  \begin{align*}
    b\circ a &= a + b + \frac{|X-Y| + |Y-Z| - |X-Z|}{2}
    \\ &= a + b + \max (x, y) + \max (y, z) - \max (x, z) - y
  \end{align*}
\end{definition}

\begin{definition}
  Fix a commutative ring $R$. The $\KLRW$ category of $1$ black strand
  and $n$ red strands, denoted $\KLRW$ category of $1$ black strand
  and $n$ red strands, denoted $\KLRW^R_n$, is given by
  \begin{itemize}
    \item $\Ob(\KLRW^R_n)$ is the set of positionings of $1$ black strand
    with $n$ red strands
    \item $\Hom(X, Y)$ is the free $R$-module generated by $\StrandSpace_{X, Y}$.
  \end{itemize}
  Composition in $\KLRW^R_n$ is defined by extending the composition of elements of the strand
  space linearly.

  Since each $\Hom(X, Y)$ is an $R$-module and composition is $R$-linear, $\KLRW^R_n$
  is $R$-linear and thus preadditive.
\end{definition}
\begin{definition}
  Let $\C$ be a preadditive category with a zero object. Then a bounded cochain
  complex of $\C$ is a ($\Z$-graded) cochain complex $A^\bullet$ of $\C$ equipped with a finite set
  $S$ called the support, satisfying
  $\forall i\in \Z, i\in S \Leftrightarrow A^i\not\cong 0$. Note that
  $S$ is uniquely determined by $A^\bullet$, though included
  for computability.
\end{definition}

\section{Bounded Cochain Complexes}
\begin{definition}
  Let $\BK(\C)$ denote the set of bounded cochain complexes in $\C$. THen
  there is a function $F:\BK(\C)\to K^\bullet(\C)$ which forgets the finiteness.
  Then we equip $\BK(\C)$ with the induced category structure of $F$ which makes
  $F$ into a fully faithful functor $F:\BK(\C) \to K^\bullet (\C)$.
\end{definition}

\section{$\beta$-functor}


\subsection{Functor data}
We will now specify the data needed to define how a braiding functor acts. Let $R, S, T \in \mathcal{KLRW}$.

The data of a braiding functor is given by the following functions:
$$
\begin{aligned}
\beta_0 &\colon & \mathcal{KLRW} &\to K^* \\
\beta_1 &\colon & \operatorname{Hom}(R,S) &\to \operatorname{Hom}(\beta_0(R), \beta_0(S)) \\
\beta_2 &\colon & \operatorname{Hom}(R,S) \times \operatorname{Hom}(S,T) &\to \operatorname{Hom}(\beta_0(R), \beta_0(T))
\end{aligned}
$$

In addition to these functions, there must also be proofs that they satisfy $A_\infty$ relations (to be specified).

One can understand these functions as specifying how the functor acts on the generating elements of our chain complex category. The action of the functor is as follows.

\subsection{Functor action}
Firstly, we may extend the functions acting on $\mathcal{KLRW}$ to act on $\operatorname{Add}(\mathcal{KLRW})$ in a linear sense. We define:
$$\widehat{\beta}_0 \colon \operatorname{Add}(\mathcal{KLRW}) \to K^*$$
by setting:
$$\widehat{\beta}_0\left(\bigoplus_i S_i\right) = \bigoplus_i \beta_0(S_i)$$

Extending the remaining two functions requires more care.

For $\beta_1$, let $\bigoplus_i A_i$ and $\bigoplus_j B_j$ be objects in $\operatorname{Add}(\mathcal{KLRW})$. Let $f \in \operatorname{Hom}\left(\bigoplus_i A_i, \bigoplus_j B_j\right)$. For each pair $k, l$, we project $f$ to its component $f_{k,l} \in \operatorname{Hom}(A_k, B_l)$. We now create $g_{k,l} \in $\operatorname{Hom}\left(\widehat{\beta}_0\left(\bigoplus_i A_i\right), \widehat{\beta}_0\left(\bigoplus_j B_j\right)\right)$ by extending this morphism. Formally, $g_{k,l}$ is constructed by composing the canonical projection, the mapped component, and the canonical injection:
$$\widehat{\beta}_0\left(\bigoplus_i A_i\right) \twoheadrightarrow {\beta}_0(A_k) \xrightarrow{\beta_1(f_{k,l})} {\beta}_0(B_l) \hookrightarrow \widehat{\beta}_0\left(\bigoplus_j B_j\right)$$



For $\beta_2$, we construct the extension $\widehat{\beta}_2$ analogously.

Note now that if $A^* \in K^*$, then we may write $A^*$ as the cone of the following chain map:
$$\begin{tikzcd}
0 \arrow[r] & A_0 \arrow[r] \arrow[d] & \dots \arrow[r] & A_{n-1} \arrow[r] \arrow[d] & 0 \\
0 \arrow[r] & 0 \arrow[r] & \dots \arrow[r] & A_n \arrow[r] & 0
\end{tikzcd}$$

We may call the upper chain complex $A^*_{[0,n-1]}$, and the lower $A^*_{[n,n]}[-1]$ (shifted by $-1$ because $A_{n-1}$ and $A_n$ need the same homological grading).

Assuming we know how $\overline{\beta}_0$ acts on these two chain complexes, we will state how it acts on $A^*$. First, take the direct sum:
$$\overline{\beta}_0(A^*) = \overline{\beta}_0(A^*_{[0,n-1]}) \oplus \overline{\beta}_0(A^*_{[n,n]}[-1])$$
